\section{Fundamentação Teórica}

\vspace{-0.4cm}

\noindent\rule{\textwidth}{0.7pt}

<Descrever sobre conceitos e fundamentos necessários ao entendimento da sua pesquisa utilizando as obras da literatura [artigos científicos, livros, teses, dissertações, tcc’s, etc], que servirão de embasamento teórico da pesquisa.

Apresentar as ideias principais dos principais autores da área. As ideias devem ser apenas apresentadas, não discutidas ou criticadas. Não devem ser incluídas as ideias do próprio autor do TCC neste capítulo.>

<É obrigatório organizar a fundamentação teórica em estrutura de tópicos.>

<Precisa que seja criado um parágrafo introdutório apresentando o capítulo de Fundamentação Teórica.>

<Os parágrafos produzidos precisam fazer uso de citações. O recomendável é 1 citação (direta ou indireta) para cada 2 parágrafos produzidos.>

\subsection{Tópico 1}

<Texto do Tópico 1>

<Os parágrafos produzidos precisam fazer uso de citações. O recomendável é 1 citação (direta ou indireta) para cada 2 parágrafos produzidos.>

\subsection{Tópico 2}

<Texto do Tópico 2>

<Os parágrafos produzidos precisam fazer uso de citações. O recomendável é 1 citação (direta ou indireta) para cada 2 parágrafos produzidos.>

\subsection{Tópico 3}

<Texto do Tópico 3>

<Os parágrafos produzidos precisam fazer uso de citações. O recomendável é 1 citação (direta ou indireta) para cada 2 parágrafos produzidos.>

\subsection{Tópico n}

<Texto do Tópico n>

<Os parágrafos produzidos precisam fazer uso de citações. O recomendável é 1 citação (direta ou indireta) para cada 2 parágrafos produzidos.>
